\section{Limitations}

The first limitation is event availability. NOAA Storm Events records are
official and hash-verified, but the public detail files do not provide
machine-readable publication timestamps. The SACB therefore uses event start as
an availability proxy, making event-conditioned results retrospective
sensitivity analyses rather than operational forecasts.

The second limitation is geographic validation. The experiment covers one
county-level PM2.5 aggregate from 2019 through 2024. It does not establish
transfer to another county, monitoring topology, climate, or emissions regime,
and therefore does not satisfy the external-validation standard for a mature
generalization claim.

The third limitation concerns model integration. The LSER uses transparent
similarity functions rather than a learned retriever, and the DFEH combines
Chronos-Bolt and retrieval only at forecast level. Event-TimeRAF therefore does
not reproduce TimeRAF's learned dual encoder or internal Channel Prompting.

The fourth limitation is the archived M12 analysis. Its source-selection log is
explicitly post hoc with respect to saved test predictions, even though the
current notebook implements validation-only selection for a future rerun. The
M12 value is consequently exploratory. Hardware model identifiers and peak
memory were also absent from the run manifest, so runtime cannot be interpreted
as a hardware-normalized efficiency result.

\section{Future Work}

The immediate priority is a clean rerun that records validation-selected M12
sources and all downstream artifacts in one manifest. Operational event studies
should then replace event-start proxies with archives that expose genuine issue
or publication times. External validation should apply the unchanged protocol
to a second US county with distinct meteorology and emissions. Longer-term work
may evaluate a learned event-context retriever, richer event embeddings,
calibrated probabilistic uncertainty, and internal TSFM conditioning. Those
extensions should be introduced only with component-wise ablations and without
changing the source audit or temporal embargo.
